% 2011 年陕西卷理科第 8 题：评分程序框图。
% 来源 PDF：普通高考/2011/2011陕西理.pdf，SHA-256 04dab6de0e3ee6223a1cbc126f7f83ae354d2b31201a14db1f31c8477d8a3bed。
% 原页：第 1 页第 2 栏上部，书页码 693。
\begin{tikzpicture}[>=stealth, semithick, font=\small, node distance=5mm]
  \node[draw, rounded corners=0.1cm, minimum width=1.3cm, minimum height=0.45cm] (start) {开始};
  \node[draw, trapezium, trapezium left angle=72, trapezium right angle=108,
        minimum width=2.4cm, minimum height=0.5cm, below=of start] (in12) {输入 $x_1$，$x_2$};
  \node[draw, diamond, aspect=2.6, minimum width=2.6cm, minimum height=0.9cm, below=of in12] (d1) {$|x_1-x_2|\leqslant2$};
  \node[draw, trapezium, trapezium left angle=72, trapezium right angle=108,
        minimum width=2.0cm, minimum height=0.5cm, below=of d1] (in3) {输入 $x_3$};
  \node[draw, diamond, aspect=3.0, minimum width=3.4cm, minimum height=1.0cm, below=of in3] (d2) {$|x_3-x_1|<|x_3-x_2|$};
  \node[draw, minimum width=1.5cm, minimum height=0.45cm, below left=8mm and 14mm of d2] (s2) {$x_2=x_3$};
  \node[draw, minimum width=1.5cm, minimum height=0.45cm, below right=8mm and 2mm of d2] (s1) {$x_1=x_3$};
  \node[draw, minimum width=2.4cm, minimum height=0.7cm, below=26mm of d2] (p) {$p=\dfrac{x_1+x_2}2$};
  \node[draw, trapezium, trapezium left angle=72, trapezium right angle=108,
        minimum width=1.8cm, minimum height=0.5cm, below=of p] (out) {输出 $p$};
  \node[draw, rounded corners=0.1cm, minimum width=1.3cm, minimum height=0.45cm, below=of out] (stop) {结束};
  \draw[->] (start)--(in12);
  \draw[->] (in12)--(d1);
  \draw[->] (d1)--node[right=2pt] {否} (in3);
  \draw[->] (in3)--(d2);
  \draw[->] (d2.west) -| node[pos=0.25, above] {是} (s2);
  \draw[->] (d2.east) -| node[pos=0.25, above] {否} (s1);
  \draw[->] (s2.south) |- (p.west);
  \draw[->] (s1.south) |- (p.east);
  \draw[->] (p)--(out);
  \draw[->] (out)--(stop);
  \draw[->] (d1.east) -- ++(2.6,0) node[pos=0.55, above] {是} |- (p.east);
\end{tikzpicture}
